% v2.3.1 figure UID: FIG-P608-01
% v2.7.0 SA2 repair: retained-sample means recomputed from the fixed trace;
% direct panel/segment labels and print-safe type/stroke hierarchy.
\tikzset{slfig-FIG-P608-01/.style={font=\fontsize{9.6pt}{11.6pt}\selectfont,
  every path/.append style={line cap=round,line join=round}}}
% Keep the two compact relation signs independently legible at the native
% 300 dpi audit scale.  The two t glyphs retain natural math script style.
\newcommand{\slfigTraceTallEq}{%
  \mathrel{\vcenter{\offinterlineskip
    \hbox{\rule{5.5pt}{.85pt}}\vskip3.9pt
    \hbox{\rule{5.5pt}{.85pt}}}}}
\newcommand{\slfigTraceScriptT}{%
  \scriptstyle t}
\begin{figure}[htbp]
\centering
\begin{tikzpicture}[slfig-FIG-P608-01]
\begin{groupplot}[group style={group size=1 by 2,vertical sep=11.0mm},
  width=11.4cm,height=3.55cm,xmin=1,xmax=20,xtick={1,5,10,15,20},
  axis lines=left,clip=false,scaled ticks=false,
  axis line style={line width=.65pt},tick style={line width=.55pt},
  tick label style={font=\fontsize{9.6pt}{11.6pt}\selectfont},
  label style={font=\fontsize{10.8pt}{13.0pt}\selectfont},
  ylabel style={rotate=-90,anchor=east,at={(axis description cs:-0.12,0.5)}},
  title style={font=\fontsize{10.8pt}{13.0pt}\selectfont,
    at={(axis description cs:0.5,1.035)},anchor=south}]
\nextgroupplot[ymin=1.4,ymax=4.05,ytick={2,3,4},ylabel={$X_{\slfigTraceScriptT}$},
  xticklabels={},title={轨迹 $X_t$}]
  \path[fill=SLSoftGray,pattern=north east lines,pattern color=SLRuleGray]
    (axis cs:1,1.4) rectangle (axis cs:5.5,4.05);
  \addplot[draw=SLBlue,line width=1.0pt,mark=*,mark size=1.7pt] coordinates
    {(1,3.8)(2,3.4)(3,3.0)(4,2.7)(5,2.4)(6,1.9)(7,2.2)(8,1.7)(9,2.1)(10,2.0)
     (11,2.3)(12,1.8)(13,2.1)(14,1.9)(15,2.2)(16,2.0)(17,1.8)(18,2.1)(19,2.0)(20,1.9)};
  \draw[draw=SLGold,line width=.85pt,densely dashed]
    (axis cs:5.5,1.4)--(axis cs:5.5,4.05);
  % Keep the warm-up annotation outside the patterned region.  A label
  % background is not used as an occlusion exemption for hatch/text clearance.
  \node[anchor=north west,font=\fontsize{9.6pt}{11.6pt}\selectfont,text=SLGray,
    align=center]
    at (axis cs:5.7,3.95) {预热段\\$t\slfigTraceTallEq 1,\ldots,5$};
  \node[anchor=north,font=\fontsize{9.6pt}{11.6pt}\selectfont,text=SLGray]
    at (axis cs:12.8,3.88) {保留样本 $t\slfigTraceTallEq 6,\ldots,20$};
\nextgroupplot[ymin=1.75,ymax=2.20,ytick={1.8,1.9,2.0,2.1},
  xlabel={$t$},ylabel={$\overline X_{6:\slfigTraceScriptT}$},
  title={保留样本运行均值 $\overline X_{6:t}$}]
  \path[fill=SLSoftGray,pattern=north east lines,pattern color=SLRuleGray]
    (axis cs:1,1.75) rectangle (axis cs:5.5,2.20);
  \addplot[draw=SLBlue,line width=1.0pt,mark=square*,mark size=1.7pt] coordinates
    {(6,1.9000)(7,2.0500)(8,1.9333)(9,1.9750)(10,1.9800)(11,2.0333)(12,2.0000)
     (13,2.0125)(14,2.0000)(15,2.0200)(16,2.0182)(17,2.0000)(18,2.0077)(19,2.0071)(20,2.0000)};
  \draw[draw=SLGold,line width=.85pt,densely dashed]
    (axis cs:5.5,1.75)--(axis cs:5.5,2.20);
  \draw[draw=SLTeal,line width=.80pt,dash dot] (axis cs:5.5,2)--(axis cs:20,2);
  \node[anchor=south east,font=\fontsize{9.6pt}{11.6pt}\selectfont,text=SLTeal]
    at (axis cs:19.6,2.135) {目标值 $2$};
\end{groupplot}
\end{tikzpicture}
\caption{舍弃前5步后，保留样本运行均值在目标值2附近波动；本图仅作诊断，不构成收敛证明}\label{fig:V5-C03-trace-running-mean}
\end{figure}
